\section{Related work}

Proof-carrying code places explicit evidence between a code producer and a
consumer's safety policy\citep{baseline01}. That relationship motivates our
separation of a proposed artifact from the evidence needed to accept it. Our
retained repair experiment uses finite executable checks; its signed records
are provenance for observations, and do not constitute proof-carrying code in
that formal sense. The conditional acceptance and reuse arguments likewise
depend on their stated premises and are not machine-checked native proofs.

Build systems provide a precise setting for reasoning about dependencies and
incremental work. Mokhov et al.\ separate the order in which tasks run from the
decision whether a task needs rebuilding\citep{baseline07}. Our design applies
that distinction to the applicability of evidence: reusing a previous test
observation additionally requires the relevant program, test, environment, and
lifecycle dependencies to remain valid. The cold A/B experiment supplies no
empirical comparison of incremental test reuse with a build system.

SWE-agent studies how a purpose-built interface for repository navigation,
editing, and program execution affects coding-agent
behavior\citep{baseline08}. This motivates treating the agent's access to
repository information as part of the method. Our intervention is the addition
of source-linked native context under a fixed proposal and scoring protocol.
Differences in tasks, interaction, model, and acceptance criteria preclude
interpreting this experiment as a head-to-head evaluation of SWE-agent or as a
new result on its benchmarks.
